% Ban thu: chi giu cac cau co trong CD2

%@ id=TT26.nguyentatthanh-d2.003 | trang=1 | so=1 | de=II | vitri=CD2.4,C3.9.5 | muc=8 | tt=0
\begin{bt}
Cho hai biểu thức $A = \dfrac{x - 4}{\sqrt{x} + 1}$ và $B = \dfrac{2}{\sqrt{x} - 2} + \dfrac{3}{\sqrt{x} + 2} + \dfrac{4\sqrt{x}}{4 - x}$. Xét tính đúng sai của các khẳng định sau:
\begin{cacy}
\item Điều kiện xác định của biểu thức $B$ là $x \ge 0$ và $x \ne 4$.
\item Giá trị của biểu thức $A$ tại $x = 4 - 2\sqrt{3}$ là $-2$.
\item Rút gọn biểu thức $B$ ta được $B = \dfrac{1}{\sqrt{x} - 2}$.
\item Cho $P = A \cdot B$. Tổng tất cả các giá trị nguyên của $x$ để $\sqrt{P} < \dfrac{1}{2}$ bằng 26.
\end{cacy}
\loigiai{
\begin{cacy}(4)
\item Đúng.
\item Đúng.
\item Sai.
\item Đúng.
\end{cacy}
}
\end{bt}
